% Acurácia do consenso ponderado por mecanismo, 65% em conluio, 20 nós.
% EMA reutiliza a série de 65% (vermelho, pontilhado, estrela): é a mesma
% curva da figura que compara os percentuais.
% Ajuste \figw e \figh abaixo.

\renewcommand{\figw}{0.90\textwidth}
\renewcommand{\figh}{7.4cm}

\begin{figure}[htbp]
\centering
\begin{tikzpicture}
\begin{axis}[tcc linha, width=\figw, height=\figh, ylabel={Acurácia}]
  \linhasDeTeste
  \addplot[serie sessenta] table[x=rodada, y=acuracia_ponderada, col sep=comma]
    {por_rodada/n20_m65_com_ema.csv};
  \addlegendentry{EMA}
  \addplot[serie ema assimetrica] table[x=rodada, y=acuracia_ponderada, col sep=comma]
    {por_rodada/n20_m65_com_ema_asymmetric.csv};
  \addlegendentry{EMA assimétrica}
  \addplot[serie beta] table[x=rodada, y=acuracia_ponderada, col sep=comma]
    {por_rodada/n20_m65_com_beta.csv};
  \addlegendentry{Beta}
  \legendaRodadaDeTeste
\end{axis}
\end{tikzpicture}
\caption{Acurácia do consenso ponderado por mecanismo de reputação.
65\% de nós maliciosos em conluio, 20 nós. Média entre as sementes.}
\label{fig:mecanismos-65}
\end{figure}
